%═══════════════════════════════════════════════════════════════════════════════
\chapter{Probabilidad, inferencia y el fundamento estadístico del aprendizaje}
\label{ch:cap7}
%═══════════════════════════════════════════════════════════════════════════════

Los seis capítulos anteriores construyeron la maquinaria técnica del aprendizaje
supervisado: la familia paramétrica, la pérdida, el optimizador, el gradiente.
Pero esa maquinaria operó con un lenguaje implícito. Se minimizó la pérdida MSE
sin preguntar por qué esa función y no otra. Se añadió regularización $\ell_2$
sin justificar su forma. Se interpretó la tasa de aprendizaje como temperatura
sin precisar de qué distribución.

Este capítulo hace explícito ese lenguaje: el lenguaje de la probabilidad y la
inferencia estadística. El resultado más importante del capítulo no es un teorema
nuevo sino una re-lectura: la pérdida MSE es el logaritmo negativo de una
verosimilitud gaussiana, la regularización $\ell_2$ es un prior gaussiano sobre
los parámetros, la distribución de Gibbs del Capítulo~\ref{ch:cap5} es la distribución de
máxima entropía bajo una restricción de energía media. Cada pieza que se construyó
pragmáticamente tiene una justificación probabilística precisa.

Este reencuadre no es cosmético. Abre tres puertas que los capítulos siguientes
atravesarán: la cuantificación de incertidumbre (qué tan seguros estamos de las
predicciones del modelo), la inferencia variacional (cómo aproximar distribuciones
intratables), y el análisis de generalización (qué garantías formales existen
sobre el comportamiento del modelo en datos nuevos). Las tres son imprescindibles
para los modelos generativos del Capítulo~\ref{ch:cap12}.

%───────────────────────────────────────────────────────────────────────────────
\section{El aprendizaje automático como inferencia estadística}
\label{sec:aprendizaje_inferencia}
\dificultad{1}{1}
%───────────────────────────────────────────────────────────────────────────────

\subsection{El modelo generativo de los datos}
\label{subsec:modelo_generativo}

Los datos $\mathcal{D} = \{(\bm{x}_i, y_i)\}_{i=1}^N$ no son fijos y
deterministas: son realizaciones de un proceso aleatorio. La hipótesis
fundamental del aprendizaje supervisado es la siguiente.

\begin{definicion}{Hipótesis de los datos i.i.d.}
\label{hip:iid}
Existe una distribución $p_{\text{data}}(\bm{x}, y)$ desconocida tal que cada
par $(\bm{x}_i, y_i)$ es una realización independiente e idénticamente
distribuida (i.i.d.) de $p_{\text{data}}$.
\end{definicion}

Esta hipótesis tiene una consecuencia inmediata: la función que se busca aproximar
no es una función arbitraria sino la \textbf{esperanza condicional}:
\[
  f^*(\bm{x}) = \mathbb{E}_{p_{\text{data}}}[y|\bm{x}]
  = \int y\,p_{\text{data}}(y|\bm{x})\,dy.
\]
$f^*$ es el predictor óptimo en el sentido del error cuadrático medio poblacional:
$f^* = \argmin_f \mathbb{E}_{p_{\text{data}}}[(f(\bm{x})-y)^2]$ (demostración
en el Ejercicio Resuelto~\ref{ej:predictor_optimo}).

\subsection{Tres fuentes de incertidumbre}
\label{subsec:tres_incertidumbres}

En cualquier problema de aprendizaje coexisten tres tipos de incertidumbre con
naturalezas distintas.

\medskip
\noindent\textbf{Incertidumbre aleatoria} (o irreducible). Es la variabilidad
intrínseca del proceso generativo: incluso conociendo $f^*$ exactamente, la
predicción $f^*(\bm{x})$ difiere de $y$ por el ruido residual $\varepsilon =
y - f^*(\bm{x})$. Esta incertidumbre no puede reducirse con más datos ni con
modelos más complejos; es una propiedad del fenómeno físico en sí. En colisiones
de partículas, la fluctuación estadística de las lluvias de partículas es
incertidumbre aleatoria.

\medskip
\noindent\textbf{Incertidumbre epistémica} (o reducible). Es la incertidumbre
sobre $\bm{\theta}$ debida a que los datos son finitos. Con $N \to \infty$,
la incertidumbre epistémica tiende a cero. Esta es la incertidumbre que la
inferencia bayesiana cuantifica mediante el posterior $p(\bm{\theta}|\mathcal{D})$,
y que el MC Dropout de la Sección~\ref{subsec:mc_dropout} estima sin inferencia
bayesiana completa.

\medskip
\noindent\textbf{Incertidumbre de modelo}. Es la incertidumbre debida a que la
familia paramétrica $\{f(\cdot;\bm{\theta})\}$ puede no contener $f^*$. Incluso
con $N \to \infty$, si la clase de funciones es demasiado restrictiva, el mejor
elemento de la familia difiere de $f^*$ por el sesgo del modelo. El Capítulo~\ref{ch:cap2}
analizó esta incertidumbre bajo el nombre de sesgo.

\begin{observacion}{La ecuación de error total}
Las tres incertidumbres contribuyen al error cuadrático medio esperado:
\[
  \mathbb{E}[(f(\bm{x};\hat{\bm{\theta}}) - y)^2]
  = \underbrace{\mathrm{Var}[\varepsilon]}_{\text{aleatoria}}
  + \underbrace{\mathrm{Sesgo}^2[f]}_{\text{de modelo}}
  + \underbrace{\mathrm{Var}[f(\bm{x};\hat{\bm{\theta}})]}_{\text{epistémica}}.
\]
Esta es exactamente la descomposición sesgo-varianza del Capítulo~\ref{ch:cap2} con una
interpretación probabilística precisa: el término de varianza es la incertidumbre
epistémica sobre los parámetros propagada a las predicciones.
\end{observacion}

%───────────────────────────────────────────────────────────────────────────────
\section{Máxima verosimilitud: el puente entre pérdida y probabilidad}
\label{sec:mle}
\dificultad{2}{2}
%───────────────────────────────────────────────────────────────────────────────

\subsection{La función de verosimilitud}
\label{subsec:verosimilitud}

\begin{definicion}{Verosimilitud y log-verosimilitud}
\label{def:verosimilitud}
Dado el modelo probabilístico $p(y|\bm{x};\bm{\theta})$, la
\textbf{verosimilitud} de los datos $\mathcal{D}$ bajo los parámetros $\bm{\theta}$
es:
\[
  \mathcal{L}_{\text{lik}}(\bm{\theta};\mathcal{D})
  = \prod_{i=1}^N p(y_i|\bm{x}_i;\bm{\theta}).
\]
La \textbf{log-verosimilitud} es:
\[
  \ell(\bm{\theta};\mathcal{D}) = \sum_{i=1}^N \log p(y_i|\bm{x}_i;\bm{\theta}).
\]
El \textbf{estimador de máxima verosimilitud} (MLE) es:
\[
  \hat{\bm{\theta}}_{\text{MLE}}
  = \argmax_{\bm{\theta}}\ell(\bm{\theta};\mathcal{D})
  = \argmin_{\bm{\theta}}\left(-\ell(\bm{\theta};\mathcal{D})\right).
\]
\end{definicion}

La product sobre las muestras refleja la Hipótesis~\ref{hip:iid}: la densidad
conjunta de muestras independientes es el producto de las densidades marginales.
El logaritmo convierte el producto en suma, facilitando el análisis y la
optimización.

\subsection{MSE como MLE bajo ruido gaussiano}
\label{subsec:mse_mle}

\begin{proposicion}{MSE como MLE}
\label{prop:mse_mle}
Si el modelo probabilístico es $p(y|\bm{x};\bm{\theta}) =
\mathcal{N}(f(\bm{x};\bm{\theta}),\sigma^2)$ con $\sigma^2$ fijo, entonces:
\[
  \hat{\bm{\theta}}_{\text{MLE}}
  = \argmin_{\bm{\theta}} \frac{1}{N}\sum_{i=1}^N (f(\bm{x}_i;\bm{\theta})-y_i)^2
  = \argmin_{\bm{\theta}} \mathcal{L}_{\text{MSE}}(\bm{\theta}).
\]
El MLE bajo ruido gaussiano aditivo e independiente es equivalente a la
minimización de la pérdida MSE.
\end{proposicion}

\begin{proof}
Con $p(y_i|\bm{x}_i;\bm{\theta}) = \frac{1}{\sqrt{2\pi\sigma^2}}
\exp\!\left(-\frac{(y_i-f(\bm{x}_i;\bm{\theta}))^2}{2\sigma^2}\right)$, la
log-verosimilitud es:
\[
  \ell(\bm{\theta}) = -\frac{N}{2}\log(2\pi\sigma^2)
  - \frac{1}{2\sigma^2}\sum_{i=1}^N (y_i-f(\bm{x}_i;\bm{\theta}))^2.
\]
Maximizar $\ell$ respecto a $\bm{\theta}$ es equivalente a minimizar
$\sum_i(y_i-f(\bm{x}_i;\bm{\theta}))^2 = N\mathcal{L}_{\text{MSE}}$.
$\square$
\end{proof}

\begin{observacion}{El principio de mínimos cuadrados de Gauss}
La Proposición~\ref{prop:mse_mle} es la justificación probabilística del método
de mínimos cuadrados, que Gauss derivó en 1809 para la estimación de órbitas
planetarias. Su argumento fue exactamente este: si los errores de medición son
gaussianos e independientes, el estimador que maximiza la probabilidad de los
datos observados es el que minimiza la suma de cuadrados. Lo que Gauss demostró
para modelos lineales en astronomía aplica, con los mismos supuestos sobre el
ruido, a cualquier familia paramétrica incluido el MLP.
\end{observacion}

\begin{fisicaconexion}{Verosimilitud y ajuste de curvas en física experimental}
El MLE bajo ruido gaussiano es la base del ajuste de curvas en física experimental.
Cuando se ajusta una señal a un modelo teórico $f(x;\bm{\theta})$ minimizando
$\chi^2 = \sum_i (y_i - f(x_i;\bm{\theta}))^2/\sigma_i^2$, se está calculando
exactamente el MLE bajo modelo gaussiano con varianzas conocidas $\sigma_i^2$.
El \textbf{estadístico $\chi^2$ reducido} $\chi^2_\nu = \chi^2/(N-p)$, con $\nu = N-p$
grados de libertad, es una prueba de bondad del ajuste: $\chi^2_\nu \approx 1$
indica que el modelo es compatible con los datos dado el nivel de ruido; $\chi^2_\nu \gg 1$
indica modelo incorrecto o ruido subestimado. Para distribuciones de Poisson
(conteos de fotones, decaimientos radiactivos, colisiones de partículas), la
log-verosimilitud es $\ell(\bm{\theta}) = \sum_i [n_i\log\lambda_i(\bm{\theta})
- \lambda_i(\bm{\theta})]$ donde $\lambda_i = f(x_i;\bm{\theta})$ es la tasa
esperada: en el límite de altas tasas converge al $\chi^2$ gaussiano, pero para
conteos bajos (condición habitual en física de altas energías) la pérdida
correcta es la log-verosimilitud de Poisson, no el MSE.
\end{fisicaconexion}

\subsection{Entropía cruzada como MLE para clasificación}
\label{subsec:ce_mle}

\begin{proposicion}{Entropía cruzada como MLE}
\label{prop:ce_mle}
Para clasificación multiclase con $K$ clases, si $p(y|\bm{x};\bm{\theta}) =
\mathrm{Categorical}(\hat{\bm{p}}(\bm{x};\bm{\theta}))$ donde
$\hat{p}_k = \mathrm{softmax}(f(\bm{x};\bm{\theta}))_k$, entonces:
\[
  \hat{\bm{\theta}}_{\text{MLE}}
  = \argmin_{\bm{\theta}} \frac{1}{N}\sum_{i=1}^N
  \left(-\sum_{k=1}^K y_{ik}\log\hat{p}_k(\bm{x}_i;\bm{\theta})\right),
\]
donde $y_{ik} = \mathbf{1}[\text{clase de }i = k]$ es la codificación one-hot.
La pérdida de entropía cruzada es el MLE bajo modelo categórico.
\end{proposicion}

\begin{proof}
La log-verosimilitud de una muestra con clase $c_i$ bajo distribución categórica es
$\log p(y_i|\bm{x}_i;\bm{\theta}) = \log\hat{p}_{c_i}(\bm{x}_i;\bm{\theta})
= \sum_k y_{ik}\log\hat{p}_k(\bm{x}_i;\bm{\theta})$. Sumando sobre las $N$
muestras y cambiando de signo (maximizar $\ell$ = minimizar $-\ell$) se obtiene
la pérdida de entropía cruzada. $\square$
\end{proof}

\subsection{La información de Fisher y la geometría del espacio de parámetros}
\label{subsec:fisher}

El MLE no solo da un estimador puntual: lleva consigo una métrica natural sobre
el espacio de parámetros que cuantifica cuánta información sobre $\bm{\theta}$
contiene una observación.

\begin{definicion}{Información de Fisher}
\label{def:fisher}
La \textbf{información de Fisher} del modelo $p(y|\bm{x};\bm{\theta})$ es la
matriz:
\[
  \bm{I}(\bm{\theta}) = \mathbb{E}_{p(y|\bm{x};\bm{\theta})}\!\left[
  \nabla_{\bm{\theta}}\log p(y|\bm{x};\bm{\theta})\,
  \nabla_{\bm{\theta}}\log p(y|\bm{x};\bm{\theta})^\top
  \right] \in \mathbb{R}^{p\times p}.
\]
\end{definicion}

La información de Fisher tiene varias interpretaciones equivalentes:

\begin{proposicion}{Propiedades de la información de Fisher}
\label{prop:fisher}
Para modelos regulares (en los que se puede intercambiar integral y derivada):
\begin{enumerate}
  \item $\bm{I}(\bm{\theta}) = -\mathbb{E}[\nabla^2_{\bm{\theta}}\log
        p(y|\bm{x};\bm{\theta})]$: es el negativo del hessiano esperado de la
        log-verosimilitud (hessiano de la log-verosimilitud negativa).
  \item La cota de Cramér-Rao: para cualquier estimador insesgado
        $\hat{\bm{\theta}}$ de $\bm{\theta}$,
        $\mathrm{Cov}(\hat{\bm{\theta}}) \succeq \bm{I}(\bm{\theta})^{-1}$
        (la covarianza del estimador es al menos la inversa de la información
        de Fisher en el sentido de matrices semidefinidas positivas).
  \item El MLE es asintóticamente eficiente: $\sqrt{N}(\hat{\bm{\theta}}_{\text{MLE}}
        - \bm{\theta}^*) \to \mathcal{N}(\bm{0}, \bm{I}(\bm{\theta}^*)^{-1})$
        en distribución.
\end{enumerate}
\end{proposicion}

La cota de Cramér-Rao establece que $\bm{I}(\bm{\theta})^{-1}$ es la varianza
mínima alcanzable para cualquier estimador insesgado: la información de Fisher
mide la capacidad máxima del modelo para aprender $\bm{\theta}$ de los datos.

\medskip
\noindent\textbf{La información de Fisher como métrica riemanniana.}
La propiedad~1 establece que $\bm{I}(\bm{\theta})$ es el hessiano esperado de la
pérdida. Esto convierte $\bm{I}(\bm{\theta})$ en una métrica natural sobre el
espacio de parámetros: la \textbf{métrica de Fisher-Rao}. En el espacio dotado
de esta métrica, la distancia entre dos distribuciones cercanas
$p(\cdot;\bm{\theta})$ y $p(\cdot;\bm{\theta}+d\bm{\theta})$ es:
\[
  ds^2 = d\bm{\theta}^\top \bm{I}(\bm{\theta})\,d\bm{\theta}.
\]
Esta métrica captura la discrepancia estadística entre distribuciones: un paso
en la dirección de mayor información requiere menos distancia paramétrica para
producir el mismo cambio en la distribución.

\medskip
\noindent\textbf{El gradiente natural.}
El descenso de gradiente ordinario minimiza $\mathcal{L}$ respecto a la métrica
euclidiana en $\mathbb{R}^p$: $\bm{\theta}^{(k+1)} = \bm{\theta}^{(k)} -
\eta\nabla\mathcal{L}$. El \textbf{descenso de gradiente natural} (Amari, 1998)
minimiza respecto a la métrica de Fisher:
\[
  \bm{\theta}^{(k+1)} = \bm{\theta}^{(k)} - \eta\,\bm{I}(\bm{\theta}^{(k)})^{-1}
  \nabla\mathcal{L}(\bm{\theta}^{(k)}).
\]
El precondicionador $\bm{I}^{-1}$ transforma el gradiente para que el paso sea
invariante bajo reparametrizaciones del modelo. Es el análogo del método de
Newton (que usa $\bm{H}^{-1}$) pero con la curvatura estadística en lugar de la
curvatura de la pérdida.

\begin{observacion}{Adam como gradiente natural diagonal}
El optimizador Adam del Capítulo~\ref{ch:cap5} usa el denominador $\sqrt{\hat{v}_j}$ para
cada parámetro, donde $\hat{v}_j \approx \mathbb{E}[(\partial\mathcal{L}/
\partial\theta_j)^2]$. Comparando con la definición de información de Fisher:
$I_{jj}(\bm{\theta}) = \mathbb{E}[(\partial\log p/\partial\theta_j)^2]$.
Para pérdidas de la forma $\mathcal{L} = -\frac{1}{N}\sum_i\log p(y_i|\bm{x}_i;
\bm{\theta})$, $\hat{v}_j \approx I_{jj}/N$. El denominador de Adam es
aproximadamente $\sqrt{I_{jj}/N}$: Adam es una aproximación diagonal al gradiente
natural. La justificación probabilística de Adam es que precondicionas el gradiente
con la información de Fisher diagonal, haciéndolo invariante a cambios de escala
en los parámetros.
\end{observacion}

%───────────────────────────────────────────────────────────────────────────────
\section{El marco bayesiano}
\label{sec:bayesiano}
\dificultad{2}{2}
%───────────────────────────────────────────────────────────────────────────────

\subsection{El teorema de Bayes para el aprendizaje}
\label{subsec:bayes_aprendizaje}

El marco bayesiano trata $\bm{\theta}$ como una variable aleatoria y actualiza
la creencia sobre $\bm{\theta}$ a la luz de los datos mediante el teorema de Bayes:
\begin{equation}
  \label{eq:bayes}
  \underbrace{p(\bm{\theta}|\mathcal{D})}_{\text{posterior}}
  = \frac{\overbrace{p(\mathcal{D}|\bm{\theta})}^{\text{verosimilitud}}\;
    \overbrace{p(\bm{\theta})}^{\text{prior}}}
    {\underbrace{p(\mathcal{D})}_{\text{evidencia}}},
\end{equation}
donde $p(\mathcal{D}) = \int p(\mathcal{D}|\bm{\theta})p(\bm{\theta})\,d\bm{\theta}$
es la constante de normalización.

El \textbf{prior} $p(\bm{\theta})$ codifica el conocimiento sobre los parámetros
antes de observar datos. El \textbf{posterior} $p(\bm{\theta}|\mathcal{D})$
actualiza ese conocimiento incorporando la información en $\mathcal{D}$.

\subsection{MAP como MLE con regularización}
\label{subsec:map}

\begin{proposicion}{MAP y regularización $\ell_2$}
\label{prop:map_l2}
El estimador de máximo a posteriori (MAP):
\[
  \hat{\bm{\theta}}_{\text{MAP}} = \argmax_{\bm{\theta}}p(\bm{\theta}|\mathcal{D})
  = \argmax_{\bm{\theta}}\bigl[\ell(\bm{\theta}) + \log p(\bm{\theta})\bigr]
\]
con prior gaussiano $p(\bm{\theta}) = \mathcal{N}(\bm{0},\tau^2\bm{I})$
es equivalente a la minimización de la pérdida MSE con regularización $\ell_2$:
\[
  \hat{\bm{\theta}}_{\text{MAP}}
  = \argmin_{\bm{\theta}}\left[\mathcal{L}_{\text{MSE}}(\bm{\theta})
  + \frac{\sigma^2}{N\tau^2}\|\bm{\theta}\|_2^2\right],
\]
con $\lambda = \sigma^2/(N\tau^2)$.
\end{proposicion}

\begin{proof}
$\log p(\bm{\theta}) = -\frac{p}{2}\log(2\pi\tau^2) - \frac{1}{2\tau^2}
\|\bm{\theta}\|_2^2$. El término constante no afecta el argmax.
Combinando con la log-verosimilitud gaussiana de la
Proposición~\ref{prop:mse_mle}:
\[
  \argmax_{\bm{\theta}}\bigl[\ell(\bm{\theta})+\log p(\bm{\theta})\bigr]
  = \argmin_{\bm{\theta}}\left[\frac{1}{2\sigma^2}\sum_i(y_i-f_i)^2
  + \frac{1}{2\tau^2}\|\bm{\theta}\|_2^2\right].
\]
Dividiendo por $N/(2\sigma^2)$ se obtiene el resultado con $\lambda =
\sigma^2/(N\tau^2)$. $\square$
\end{proof}

El círculo se cierra: la regularización de Tikhonov del Capítulo~\ref{ch:cap2}, la
penalización $\ell_2$ de los Capítulos~2 y~3, y el estimador MAP bajo prior
gaussiano son el mismo objeto visto desde tres perspectivas distintas. El
coeficiente de regularización $\lambda$ es el cociente entre la varianza del ruido
y la varianza del prior: expresa cuánto peso se le da a los datos vs.\ al
conocimiento previo.

\subsection{La predicción bayesiana completa}
\label{subsec:prediccion_bayesiana}

El MLE y el MAP dan estimadores puntuales de $\bm{\theta}$. El enfoque bayesiano
completo integra sobre todos los parámetros ponderados por el posterior:
\begin{equation}
  \label{eq:pred_bayesiana}
  p(y^*|\bm{x}^*, \mathcal{D})
  = \int p(y^*|\bm{x}^*,\bm{\theta})\,p(\bm{\theta}|\mathcal{D})\,d\bm{\theta}.
\end{equation}
La predicción bayesiana completa entrega una distribución sobre $y^*$, no solo
un valor: cuantifica automáticamente la incertidumbre epistémica sobre la
predicción.

Para el MLP con $p \sim 10^6$ parámetros, la integral~\eqref{eq:pred_bayesiana}
es intractable: no existe forma analítica y el muestreo de Monte Carlo puro en
$\mathbb{R}^{10^6}$ es computacionalmente imposible. Esto motiva la inferencia
variacional de la Sección~\ref{sec:vi}.

\subsection{Priors y el principio de máxima entropía}
\label{subsec:maxent}

La elección del prior no es arbitraria. El \textbf{principio de máxima entropía}
(Jaynes, 1957) establece que el prior correcto bajo una restricción dada es la
distribución que maximiza la entropía sujeta a esa restricción: es la distribución
"menos informativa" compatible con el conocimiento disponible.

\begin{proposicion}{Distribuciones de máxima entropía}
\label{prop:maxent}
\begin{enumerate}
  \item Bajo la restricción $\mathbb{E}[\|\bm{\theta}\|_2^2] \leq C$ (norma
        cuadrática acotada), la distribución de máxima entropía es el prior
        gaussiano $\mathcal{N}(\bm{0},\sigma^2\bm{I})$.
  \item Bajo la restricción $\mathbb{E}[\mathcal{L}(\bm{\theta})] = U$ (energía
        media fija), la distribución de máxima entropía es la distribución de
        Gibbs $p^*(\bm{\theta}) \propto e^{-\mathcal{L}(\bm{\theta})/T}$.
\end{enumerate}
\end{proposicion}

La proposición~2 conecta directamente con el Capítulo~\ref{ch:cap5}: la distribución
estacionaria del SGD, $p^* \propto e^{-2\mathcal{L}/(\eta c)}$, es la
distribución de máxima entropía bajo la restricción de que la pérdida esperada
sea igual al valor que fija el ruido del mini-lote. La temperatura $T = \eta c/2$
es el multiplicador de Lagrange de esa restricción. Sin necesidad de mecánica
estadística formal, el principio de máxima entropía justifica por qué el SGD
converge a esa distribución y por qué la temperatura tiene el significado que se
describió.

%───────────────────────────────────────────────────────────────────────────────
\section{Entropía, divergencia KL e información}
\label{sec:entropia_kl}
\dificultad{2}{2}
%───────────────────────────────────────────────────────────────────────────────

\subsection{Entropía de Shannon}
\label{subsec:entropia_shannon}

\begin{definicion}{Entropía de Shannon}
\label{def:entropia}
La \textbf{entropía} de una distribución discreta $p = (p_1,\ldots,p_K)$ es:
\[
  H(p) = -\sum_{k=1}^K p_k\log p_k \geq 0,
\]
con la convención $0\log 0 = 0$. Para una distribución continua con densidad
$p(\bm{x})$, la \textbf{entropía diferencial} es:
\[
  h(p) = -\int p(\bm{x})\log p(\bm{x})\,d\bm{x}.
\]
\end{definicion}

La entropía mide la \textbf{incertidumbre} de la distribución. Sus propiedades
fundamentales:
\begin{enumerate}
  \item $H(p) \geq 0$, con igualdad si y solo si $p$ es un punto masa (certeza
        completa).
  \item $H(p) \leq \log K$, con igualdad para la distribución uniforme
        (máxima incertidumbre).
  \item Aditividad: para variables independientes $X, Y$,
        $H(X,Y) = H(X) + H(Y)$.
  \item Interpretación de teoría de la información: $H(p)$ es el número mínimo
        de bits necesarios en promedio para codificar un mensaje de la
        distribución $p$ (teorema de Shannon, 1948).
\end{enumerate}

\begin{fisicaconexion}{Entropía de Boltzmann y entropía de Shannon}
La entropía de Shannon $H(p) = -\sum_k p_k \log p_k$ y la entropía estadística
de Boltzmann $S = -k_B\sum_i p_i\ln p_i$ son el mismo objeto: la diferencia es
únicamente la constante $k_B$ (constante de Boltzmann) y la base del logaritmo.
Boltzmann derivó $S$ en 1877 como medida del número de microestados compatibles
con un macroestado; Shannon la redescubrió en 1948 como medida de la información.
La \textbf{entropía máxima} de la mecánica estadística ---que la distribución de
equilibrio es la que maximiza $S$ bajo la restricción de energía media fija---
es exactamente la Proposición~\ref{prop:maxent}: la distribución de Gibbs como
distribución de máxima entropía. La temperatura termodinámica $T$ es el
multiplicador de Lagrange de esa restricción, que en el SGD se identifica con
$\eta c/2$ (Capítulo~\ref{ch:cap5}). La \textbf{tercera ley de la termodinámica}
($S \to 0$ cuando $T \to 0$) tiene su análogo en el aprendizaje: cuando $\eta \to 0$
la distribución estacionaria del SGD se concentra en los mínimos de
$\mathcal{L}$, y la entropía de esa distribución tiende a cero.
\end{fisicaconexion}

\subsection{Divergencia de Kullback-Leibler}
\label{subsec:kl}

\begin{definicion}{Divergencia de Kullback-Leibler}
\label{def:kl}
La \textbf{divergencia KL} de $q$ a $p$ es:
\[
  D_{\text{KL}}(p\|q)
  = \int p(\bm{x})\log\frac{p(\bm{x})}{q(\bm{x})}\,d\bm{x}
  = \mathbb{E}_p\!\left[\log\frac{p(\bm{x})}{q(\bm{x})}\right].
\]
\end{definicion}

\begin{proposicion}{No negatividad de la KL}
\label{prop:kl_nn}
$D_{\text{KL}}(p\|q) \geq 0$ para todo par de distribuciones $p,q$, con igualdad
si y solo si $p = q$ casi en todo punto.
\end{proposicion}

\begin{proof}
$-\log$ es estrictamente convexa. Por la desigualdad de Jensen:
\[
  D_{\text{KL}}(p\|q)
  = \mathbb{E}_p\!\left[-\log\frac{q}{p}\right]
  \geq -\log\mathbb{E}_p\!\left[\frac{q}{p}\right]
  = -\log\int q(\bm{x})\,d\bm{x} = -\log 1 = 0.
\]
La igualdad se da cuando $\frac{q}{p}$ es constante, es decir $q = p$.
$\square$
\end{proof}

\begin{observacion}{La KL no es una distancia}
$D_{\text{KL}}(p\|q) 
eq D_{\text{KL}}(q\|p)$ en general: la KL es asimétrica.
Tampoco satisface la desigualdad triangular. No es una distancia en sentido
métrico. Sin embargo, la desigualdad de Pinsker vincula la KL con la distancia de
variación total $\mathrm{TV}(p,q) = \frac{1}{2}\int|p-q|\,d\bm{x}$:
$\mathrm{TV}(p,q)^2 \leq \frac{1}{2}D_{\text{KL}}(p\|q)$.
\end{observacion}

\subsection{Las dos direcciones de la KL y sus consecuencias}
\label{subsec:kl_direcciones}

La asimetría de la KL tiene consecuencias prácticas profundas.

\medskip
\noindent\textbf{KL forward $D_{\text{KL}}(p\|q)$.} Penaliza regiones donde $p$
tiene masa pero $q$ no: si $p(\bm{x}) > 0$ y $q(\bm{x}) \approx 0$, el término
$p\log(p/q) \to \infty$. Minimizar $D_{\text{KL}}(p\|q)$ respecto a $q$ fuerza
a $q$ a cubrir todo el soporte de $p$. En inferencia variacional (Sección~\ref{sec:vi}),
se minimiza la KL reversa $D_{\text{KL}}(q\|p)$ (más manejable), no la forward.

\medskip
\noindent\textbf{KL reversa $D_{\text{KL}}(q\|p)$.} Penaliza regiones donde
$q$ tiene masa pero $p$ no: fuerza a $q$ a concentrarse en los modos de $p$.
Para distribuciones multimodales, minimizar $D_{\text{KL}}(q\|p)$ con $q$
unimodal produce una $q$ que cubre solo uno de los modos.

\subsection{Entropía cruzada y su relación con la KL}
\label{subsec:ce_kl}

\begin{proposicion}{Descomposición de la entropía cruzada}
\label{prop:ce_descomp}
La entropía cruzada entre $p$ y $q$ satisface:
\[
  H(p,q) = -\int p(\bm{x})\log q(\bm{x})\,d\bm{x}
  = H(p) + D_{\text{KL}}(p\|q).
\]
\end{proposicion}

\begin{proof}
$H(p,q) = -\int p\log q = -\int p\log p + \int p\log(p/q)
= H(p) + D_{\text{KL}}(p\|q)$. $\square$
\end{proof}

La Proposición~\ref{prop:ce_descomp} da la interpretación precisa de la pérdida
de entropía cruzada: minimizar $H(p_{\text{data}}, p_{\bm{\theta}})$ respecto a
$\bm{\theta}$ equivale a minimizar $D_{\text{KL}}(p_{\text{data}}\|p_{\bm{\theta}})$
(la entropía de los datos $H(p_{\text{data}})$ no depende de $\bm{\theta}$).
Entrenar con pérdida de entropía cruzada es minimizar la divergencia KL entre
la distribución del modelo y la distribución de los datos: el entrenamiento busca
que el modelo sea estadísticamente indistinguible de los datos.

\subsection{Información mutua}
\label{subsec:info_mutua}

\begin{definicion}{Información mutua}
\label{def:info_mutua}
La \textbf{información mutua} entre variables aleatorias $X$ e $Y$ es:
\[
  I(X;Y) = D_{\text{KL}}\bigl(p(x,y)\|p(x)p(y)\bigr)
  = H(X) - H(X|Y) = H(Y) - H(Y|X).
\]
\end{definicion}

$I(X;Y) \geq 0$, con igualdad si y solo si $X$ e $Y$ son independientes.
Mide cuánta información sobre $X$ contiene $Y$: cuánto se reduce la incertidumbre
sobre $X$ al conocer $Y$.

En redes neuronales, la información mutua $I(\bm{a}^{(l)};\bm{y})$ entre la
activación de la capa $l$ y la etiqueta cuantifica cuánta información relevante
ha extraído esa capa. El principio del \textbf{cuello de botella de información}
(\textit{information bottleneck}, Tishby et al., 2000) propone que el aprendizaje
consiste en comprimir la representación intermedia para retener la información
relevante sobre $\bm{y}$ y descartar la irrelevante sobre $\bm{x}$: maximizar
$I(\bm{a}^{(l)};\bm{y})$ mientras se minimiza $I(\bm{a}^{(l)};\bm{x})$.

%───────────────────────────────────────────────────────────────────────────────
\section{Inferencia variacional}
\label{sec:vi}
\dificultad{3}{2}
%───────────────────────────────────────────────────────────────────────────────

\subsection{El problema de la integración sobre el posterior}
\label{subsec:prob_integracion}

La predicción bayesiana~\eqref{eq:pred_bayesiana} requiere integrar sobre el
posterior. Para el MLP con $p \sim 10^6$, el posterior $p(\bm{\theta}|\mathcal{D})$
vive en $\mathbb{R}^{10^6}$ y no tiene forma analítica. La integración de Monte
Carlo directa requiere muestras del posterior, lo que a su vez requiere MCMC en
alta dimensión: inviable para el entrenamiento iterativo.

La \textbf{inferencia variacional} convierte el problema de integración en un
problema de optimización: en lugar de muestrear del posterior, se busca la
distribución más cercana en una familia tractable.

\subsection{La ELBO: evidencia como cota inferior}
\label{subsec:elbo}

\begin{definicion}{La ELBO}
\label{def:elbo}
Dada una familia variacional $\mathcal{Q} = \{q_{\bm{\phi}}(\bm{\theta})\}$,
la \textbf{cota inferior de la evidencia} (ELBO, \textit{Evidence Lower BOund})
es:
\[
  \mathrm{ELBO}(q_{\bm{\phi}})
  = \mathbb{E}_{q_{\bm{\phi}}}\bigl[\log p(\mathcal{D}|\bm{\theta})\bigr]
  - D_{\text{KL}}\bigl(q_{\bm{\phi}}(\bm{\theta})\|p(\bm{\theta})\bigr).
\]
\end{definicion}

\begin{teorema}{Descomposición de la evidencia}
\label{thm:elbo_descomp}
Para cualquier $q_{\bm{\phi}}$:
\[
  \log p(\mathcal{D}) = \mathrm{ELBO}(q_{\bm{\phi}})
  + D_{\text{KL}}\bigl(q_{\bm{\phi}}\|p(\bm{\theta}|\mathcal{D})\bigr).
\]
Como $D_{\text{KL}} \geq 0$, la ELBO es una cota inferior de $\log p(\mathcal{D})$.
Maximizar la ELBO respecto a $\bm{\phi}$ minimiza $D_{\text{KL}}(q_{\bm{\phi}}\|
p(\bm{\theta}|\mathcal{D}))$.
\end{teorema}

\begin{proof}
\begin{align*}
  \mathrm{ELBO}(q_{\bm{\phi}})
  &= \mathbb{E}_{q_{\bm{\phi}}}[\log p(\mathcal{D}|\bm{\theta})]
    - D_{\text{KL}}(q_{\bm{\phi}}\|p(\bm{\theta})) \\
  &= \mathbb{E}_{q_{\bm{\phi}}}\!\left[\log p(\mathcal{D}|\bm{\theta})
    + \log p(\bm{\theta}) - \log q_{\bm{\phi}}(\bm{\theta})\right] \\
  &= \mathbb{E}_{q_{\bm{\phi}}}\!\left[\log\frac{p(\bm{\theta}|\mathcal{D})
    p(\mathcal{D})}{q_{\bm{\phi}}(\bm{\theta})}\right] \\
  &= \log p(\mathcal{D})
    - D_{\text{KL}}(q_{\bm{\phi}}\|p(\bm{\theta}|\mathcal{D})).
\end{align*}
Reordenando: $\log p(\mathcal{D}) = \mathrm{ELBO} + D_{\text{KL}}(q_{\bm{\phi}}\|
p(\cdot|\mathcal{D}))$. $\square$
\end{proof}

La ELBO tiene una interpretación dual que el físico reconocerá:
\[
  \mathrm{ELBO}
  = \underbrace{\mathbb{E}_{q_{\bm{\phi}}}[\log p(\mathcal{D}|\bm{\theta})]}_{\text{reconstrucción (verosimilitud esperada)}}
  - \underbrace{D_{\text{KL}}(q_{\bm{\phi}}\|p(\bm{\theta}))}_{\text{regularización (complejidad del modelo)}}.
\]
El primer término mide qué tan bien el modelo (bajo la distribución variacional
$q$) explica los datos. El segundo penaliza cuánto se aparta $q$ del prior:
es un término de regularización bayesiana. La ELBO equilibra ajuste y complejidad.

\begin{fisicaconexion}{La ELBO como energía libre de Helmholtz}
La ELBO tiene una interpretación termodinámica exacta. Escribiendo
$-\mathrm{ELBO} = \mathbb{E}_q[-\log p(\mathcal{D}|\bm{\theta})] + D_\mathrm{KL}(q\|p)$,
se identifican los dos términos con la \textbf{energía libre de Helmholtz}
$F = U - TS$:
\[
  -\mathrm{ELBO} = \underbrace{\mathbb{E}_q[-\log p(\mathcal{D}|\bm{\theta})]}_{\text{energía media } U}
  - \underbrace{H(q)}_{\text{entropía de la distribución variacional } TS}.
\]
La energía media $U = \mathbb{E}_q[-\log p(\mathcal{D}|\bm{\theta})]$ mide el
ajuste a los datos; la entropía $H(q)$ mide cuánto ``ocupa'' la distribución
variacional el espacio de parámetros. Maximizar la ELBO es minimizar $F = U - TS$:
exactamente el principio de mínima energía libre que rige el equilibrio
termodinámico. Cuanto mayor la temperatura (mayor entropía), más dispersa es
$q$ sobre el espacio de parámetros (menos sobreajuste); cuanto menor, más
concentrada en el MAP (mayor ajuste a los datos). El término de regularización
$D_\mathrm{KL}(q\|p)$ mide la divergencia del prior, análogo a la energía libre
de exceso respecto al estado de referencia. Esta analogía no es accidental: la
inferencia variacional y la física estadística son formalismos duales del mismo
principio de optimización bajo incertidumbre.
\end{fisicaconexion}

\subsection{El truco de reparametrización}
\label{subsec:reparam}

Para maximizar la ELBO con descenso de gradiente se necesita:
\[
  \nabla_{\bm{\phi}}\mathbb{E}_{q_{\bm{\phi}}(\bm{\theta})}[f(\bm{\theta})].
\]
El problema es que la esperanza depende de $\bm{\phi}$ a través de la distribución
de muestreo, y el gradiente no puede pasarse dentro de la esperanza directamente
cuando el dominio de integración depende de $\bm{\phi}$.

\begin{definicion}{Truco de reparametrización}
\label{def:reparam}
Si existe una función diferenciable $g_{\bm{\phi}}$ y una distribución base
$p(\bm{\varepsilon})$ independiente de $\bm{\phi}$ tales que
$\bm{\theta} = g_{\bm{\phi}}(\bm{\varepsilon})$ con
$\bm{\varepsilon} \sim p(\bm{\varepsilon})$ tiene la misma distribución que
$q_{\bm{\phi}}$, entonces:
\[
  \nabla_{\bm{\phi}}\mathbb{E}_{q_{\bm{\phi}}}[f(\bm{\theta})]
  = \mathbb{E}_{p(\bm{\varepsilon})}\!\left[
    \nabla_{\bm{\phi}} f(g_{\bm{\phi}}(\bm{\varepsilon}))
  \right].
\]
\end{definicion}

Para la familia gaussiana diagonal $q_{\bm{\phi}} = \mathcal{N}(\bm{\mu},
\mathrm{diag}(\bm{\sigma}^2))$ con $\bm{\phi} = (\bm{\mu},\log\bm{\sigma})$,
la reparametrización es:
\[
  \bm{\theta} = \bm{\mu} + \bm{\sigma}\odot\bm{\varepsilon},
  \qquad \bm{\varepsilon} \sim \mathcal{N}(\bm{0},\bm{I}).
\]
El gradiente de la ELBO respecto a $\bm{\phi}$ se convierte en:
\[
  \nabla_{\bm{\phi}}\,\mathrm{ELBO}
  \approx \frac{1}{K}\sum_{k=1}^K \nabla_{\bm{\phi}}
  \left[\log p(\mathcal{D}|g_{\bm{\phi}}(\bm{\varepsilon}^{(k)}))
  - \log\frac{q_{\bm{\phi}}(g_{\bm{\phi}}(\bm{\varepsilon}^{(k)}))}{p(g_{\bm{\phi}}(\bm{\varepsilon}^{(k)}))}\right],
\]
estimado con $K$ muestras $\bm{\varepsilon}^{(k)} \sim \mathcal{N}(\bm{0},\bm{I})$.
Con $K=1$ esto es el estimador SGVB (\textit{Stochastic Gradient Variational Bayes})
de Kingma y Welling (2013): un solo forward pass por iteración, compatible con
el ciclo de entrenamiento estándar del Capítulo~\ref{ch:cap5}.

%───────────────────────────────────────────────────────────────────────────────
\section{Redes neuronales bayesianas y cuantificación de incertidumbre}
\label{sec:bnn}
\dificultad{3}{2}
%───────────────────────────────────────────────────────────────────────────────

\subsection{La red neuronal bayesiana}
\label{subsec:bnn}

\begin{definicion}{Red neuronal bayesiana}
\label{def:bnn}
Una \textbf{red neuronal bayesiana} (BNN) es un MLP donde los parámetros
$\bm{\theta}$ se tratan como variables aleatorias con prior $p(\bm{\theta})$.
La predicción es la distribución predictiva posterior:
\[
  p(y^*|\bm{x}^*,\mathcal{D})
  = \int p(y^*|\bm{x}^*,\bm{\theta})\,p(\bm{\theta}|\mathcal{D})\,d\bm{\theta}.
\]
\end{definicion}

Con la familia variacional gaussiana diagonal y el truco de reparametrización,
la ELBO se maximiza con backpropagation estándar. Los parámetros variacionales
son $\bm{\phi} = (\bm{\mu},\log\bm{\sigma}) \in \mathbb{R}^{2p}$: el doble
de los parámetros de una red determinista. En lugar de aprender un solo valor
para cada peso, la BNN aprende la media y el logaritmo de la desviación estándar
de cada peso.

\subsection{Dropout como inferencia variacional aproximada}
\label{subsec:dropout_vi}

El resultado de Gal y Ghahramani (2016) establece que el entrenamiento con
dropout es una forma de inferencia variacional aproximada.

\begin{proposicion}{Dropout como inferencia variacional}
\label{prop:dropout_vi}
Considérese una BNN con prior $p(\bm{W}^{(l)}) = \prod_j\mathcal{N}(\bm{w}_j;
\bm{0},\bm{I}/l^2)$ y familia variacional formada por matrices con columnas
$q(\bm{w}_j) = p_d\,\delta(\bm{w}_j) + (1-p_d)\,\mathcal{N}(\bm{w}_j;
\bm{m}_j,\sigma^2\bm{I})$ (mezcla de punto en cero y gaussiana). La ELBO
de esta BNN, aproximada ignorando términos de varianza, es proporcional a la
pérdida de entrenamiento con dropout con probabilidad $p_d$ más un término de
regularización en las medias $\bm{m}_j$:
\[
  -\mathrm{ELBO}(q)
  \propto \mathcal{L}_{\text{dropout}}(\hat{\bm{M}})
  + \frac{p_d}{2l^2}\sum_{l,j}\|\bm{m}_j^{(l)}\|_2^2
  + \mathrm{const.}
\]
\end{proposicion}

La demostración completa se desarrolla en el Ejercicio Resuelto~\ref{ej:dropout_vi}.
La consecuencia práctica es el \textbf{MC Dropout}.

\begin{definicion}{MC Dropout para cuantificación de incertidumbre}
\label{def:mc_dropout}
\label{subsec:mc_dropout}
Para una red entrenada con dropout, la distribución predictiva posterior se
aproxima haciendo $K$ evaluaciones con dropout activo:
\begin{align*}
  \mathbb{E}[y^*|\bm{x}^*,\mathcal{D}]
  &\approx \frac{1}{K}\sum_{k=1}^K f(\bm{x}^*;\hat{\bm{\theta}}_k), \\
  \mathrm{Var}[y^*|\bm{x}^*,\mathcal{D}]
  &\approx \sigma^2 + \frac{1}{K}\sum_{k=1}^K
    \bigl(f(\bm{x}^*;\hat{\bm{\theta}}_k)
    - \mathbb{E}[y^*]\bigr)^2,
\end{align*}
donde $\hat{\bm{\theta}}_k$ es el conjunto de pesos activos en la $k$-ésima
evaluación con dropout.
\end{definicion}

El primer término de la varianza ($\sigma^2$) es la incertidumbre aleatoria
(ruido del modelo); el segundo es la incertidumbre epistémica (variabilidad
entre las distintas subred). MC Dropout cuantifica ambas sin el costo
computacional adicional de entrenar una BNN completa: se usa el mismo modelo
entrenado con dropout, simplemente activando dropout en test.

\subsection{Calibración de modelos}
\label{subsec:calibracion}

Un modelo bien calibrado satisface: cuando predice probabilidad $\hat{p}$ para
un evento, ese evento ocurre con frecuencia $\hat{p}$. Formalmente, para
clasificación binaria:
\[
  P(y=1 \mid \hat{p}(\bm{x}) = p) = p, \qquad \forall\, p \in [0,1].
\]

\begin{definicion}{Error de calibración esperado}
\label{def:ece}
El \textbf{error de calibración esperado} (ECE) se calcula dividiendo el rango
$[0,1]$ en $M$ bins $B_m$ e igualando la confianza media con la precisión media:
\[
  \mathrm{ECE} = \sum_{m=1}^M \frac{|B_m|}{N}
  \left|\mathrm{acc}(B_m) - \mathrm{conf}(B_m)\right|,
\]
donde $\mathrm{acc}(B_m)$ es la precisión en el bin y $\mathrm{conf}(B_m)$
la confianza media.
\end{definicion}

Los MLPs entrenados con entropía cruzada tienden a estar \textbf{sobreconfiados}:
predicen probabilidades cercanas a 0 o 1 con más frecuencia de lo que su precisión
justifica. La calibración puede corregirse con \textit{temperature scaling}
(dividir los logits por una constante $T > 1$ antes del softmax) o con la
regresión de Platt. En física de partículas, donde las predicciones del modelo
se usan para tomar decisiones sobre la existencia de nueva física, la calibración
correcta es tan importante como la exactitud media.

%───────────────────────────────────────────────────────────────────────────────
\section{Cotas de generalización y la teoría del aprendizaje estadístico}
\label{sec:generalizacion}
\dificultad{3}{2}
%───────────────────────────────────────────────────────────────────────────────

\subsection{El error de generalización y la desigualdad de Hoeffding}
\label{subsec:hoeffding}

El objetivo del aprendizaje es minimizar el \textbf{error poblacional} (o de
generalización):
\[
  E_{\text{pop}}(\bm{\theta}) = \mathbb{E}_{p_{\text{data}}}[\ell(f(\bm{x};\bm{\theta}),y)],
\]
que no es observable. Se minimiza en su lugar el \textbf{error empírico}:
\[
  E_{\text{train}}(\bm{\theta}) = \frac{1}{N}\sum_{i=1}^N\ell(f(\bm{x}_i;\bm{\theta}),y_i).
\]
La brecha entre ambos es el \textbf{error de generalización}:
$E_{\text{gen}} = E_{\text{pop}} - E_{\text{train}} \geq 0$.

Para un único $\bm{\theta}$ fijo (no elegido en función de los datos), la
desigualdad de Hoeffding acota esta brecha.

\begin{teorema}{Desigualdad de Hoeffding}
\label{thm:hoeffding}
Sean $X_1,\ldots,X_N$ variables aleatorias i.i.d. con $a \leq X_i \leq b$ y
$\mathbb{E}[X_i] = \mu$. Para todo $t > 0$:
\[
  P\!\left(\left|\frac{1}{N}\sum_{i=1}^N X_i - \mu\right| \geq t\right)
  \leq 2\exp\!\left(-\frac{2N t^2}{(b-a)^2}\right).
\]
\end{teorema}

Aplicando Hoeffding a la pérdida (asumida acotada en $[0,1]$ por simplicidad):
con probabilidad $\geq 1-\delta$ sobre la muestra de entrenamiento:
\[
  E_{\text{pop}}(\bm{\theta}) \leq E_{\text{train}}(\bm{\theta})
  + \sqrt{\frac{\ln(2/\delta)}{2N}}.
\]
Esta cota es razonable para un $\bm{\theta}$ fijo, pero el aprendizaje elige
$\bm{\theta}$ en función de los datos: hay que controlar el máximo sobre todos
los $\bm{\theta}$ de la familia.

\subsection{Dimensión de Vapnik-Chervonenkis y cotas de capacidad}
\label{subsec:vc}

\begin{definicion}{VC-dimensión}
\label{def:vc}
La \textbf{dimensión de Vapnik-Chervonenkis} (VC-dimensión) de una clase de
funciones binarias $\mathcal{F}$ es el tamaño máximo de un conjunto de puntos
que $\mathcal{F}$ puede \textbf{fragmentar} (\textit{shatter}): clasificar de
todas las $2^d$ formas posibles. Formalmente:
\[
  \mathrm{VCdim}(\mathcal{F}) = \max\{d : \exists\,\bm{x}_1,\ldots,\bm{x}_d
  \text{ s.t. } \mathcal{F} \text{ fragmenta } \{\bm{x}_1,\ldots,\bm{x}_d\}\}.
\]
\end{definicion}

\begin{teorema}{Cota de generalización por VC (Vapnik-Chervonenkis, 1971)}
\label{thm:vc_bound}
Sea $\mathcal{F}$ una clase de clasificadores binarios con $\mathrm{VCdim}
(\mathcal{F}) = d < \infty$. Con probabilidad $\geq 1-\delta$ sobre la muestra
de entrenamiento de tamaño $N$:
\[
  \sup_{\hat{f}\in\mathcal{F}}\bigl[E_{\text{pop}}(\hat{f})
  - E_{\text{train}}(\hat{f})\bigr]
  \leq \sqrt{\frac{8d\ln(2N/d) + 8\ln(4/\delta)}{N}}.
\]
\end{teorema}

Para MLPs la VC-dimensión escala como $O(p\log p)$ donde $p$ es el número de
parámetros (Bartlett, 1998), lo que da una cota de generalización de orden
$O(\sqrt{p\log p/N})$.

\subsection{Las limitaciones de las cotas de VC y el \textit{double descent}}
\label{subsec:limitaciones_vc}

Las cotas de VC son el resultado central de la teoría clásica del aprendizaje,
pero tienen limitaciones importantes que deben comprenderse para no usarlas mal.

\medskip
\noindent\textbf{Limitación 1: las cotas son genéricamente flojas.}
La cota del Teorema~\ref{thm:vc_bound} es del orden $O(\sqrt{p/N})$: crece al
aumentar $p$. Para un MLP con $p = 10^6$ y $N = 10^4$ daría una cota de
generalización $\gg 1$, lo que es trivialmente inútil (cualquier error es
$\leq 1$ para clasificación binaria). En la práctica, esas redes generalizan
bien. Las cotas de VC no explican por qué.

\medskip
\noindent\textbf{Limitación 2: ignoran la estructura del algoritmo.}
Las cotas de VC controlan el peor caso sobre toda la clase $\mathcal{F}$. No
capturan el hecho de que el SGD con cierta inicialización y tasa de aprendizaje
no explora toda la clase: tiene un sesgo implícito hacia soluciones con propiedades
específicas (mínima norma, mínimos planos). La regularización implícita del
Capítulo~\ref{ch:cap8} hace que el modelo efectivo tenga una complejidad mucho menor que la
VC-dimensión nominal.

\medskip
\noindent\textbf{Limitación 3: el \textit{double descent} contradice la predicción clásica.}
La cota del Teorema~\ref{thm:vc_bound} implica que el error de generalización
crece monótonamente con $p$ (para $N$ fijo). El fenómeno \textit{double descent}
(Capítulo~\ref{ch:cap8}) muestra que para $p \gg N$, el error de generalización puede volver
a decrecer. Las cotas de VC no tienen mecanismo para explicar este comportamiento:
son intrínsecamente cotas del régimen subparamétrico ($p < N$).

\medskip
\noindent\textbf{La cota de VC como herramienta de comprensión, no de predicción.}
La VC-dimensión es útil para entender cualitativamente qué familias son más o
menos complejas, para demostrar que el aprendizaje es posible en principio (el
teorema fundamental del PAC-learning: si $\mathrm{VCdim}(\mathcal{F}) < \infty$,
la familia es PAC-aprendible), y para diseñar algoritmos con garantías en el
régimen de datos abundantes ($N \gg p$). No es una herramienta predictiva en el
régimen sobredeterminado moderno.

\begin{observacion}{Cotas modernas: la teoría del margen y la norma}
La investigación reciente (Bartlett et al., 2020; Neyshabur et al., 2018)
muestra que las cotas de generalización para redes neuronales deben controlar
no la VC-dimensión sino la \textbf{norma de los parámetros} (norma espectral
o norma de Frobenius de los pesos). Estas cotas escalan como $O(\|\bm{W}\|/
\sqrt{N})$ y pueden ser significativas incluso para $p \gg N$ si los pesos
son pequeños. Son las cotas correctas para el régimen moderno, pero requieren
herramientas de teoría de matrices aleatorias para su análisis
en el límite termodinámico $p,N\to\infty$ con $p/N\to\gamma$ constante.
\end{observacion}

%───────────────────────────────────────────────────────────────────────────────
\section{Resumen del Capítulo 7}
\label{sec:resumen7}
%───────────────────────────────────────────────────────────────────────────────

Este capítulo proveyó el lenguaje probabilístico que los capítulos anteriores
usaron de forma implícita y que los capítulos siguientes necesitarán
explícitamente.

El hilo central fue la unificación: la pérdida MSE es MLE gaussiano, la
regularización $\ell_2$ es un prior gaussiano bajo MAP, la distribución
estacionaria del SGD es la distribución de máxima entropía bajo restricción
de energía. El MLE, el MAP y la inferencia variacional son tres puntos en un
continuo que va del estimador puntual a la distribución completa sobre los
parámetros. La entropía cruzada minimiza la KL entre el modelo y los datos;
la ELBO maximiza una cota inferior de la evidencia; el truco de reparametrización
hace que ambos sean diferenciables y compatibles con backpropagation.

Las cotas de VC formalizan la intuición del Capítulo~\ref{ch:cap2} sobre el compromiso
sesgo-varianza, pero tienen limitaciones fundamentales en el régimen sobredeterminado
moderno. Esas limitaciones señalan hacia dos desarrollos que el libro abordará
más adelante: la regularización implícita del SGD (Capítulo~\ref{ch:cap8}) y el análisis de
matrices aleatorias en el límite termodinámico.

Tres conexiones hacia adelante. El Capítulo~\ref{ch:cap8} usará la distinción entre
incertidumbre epistémica y aleatoria para el MC Dropout, y la ELBO para el
análisis del fenómeno \textit{double descent} en términos de complejidad efectiva.
El Capítulo~\ref{ch:cap10} completará el cálculo de Itô con rigor formal y mostrará que la
SDE de Langevin del SGD es el proceso cuya distribución estacionaria es exactamente
la distribución de Gibbs derivada aquí como distribución de máxima entropía. El
Capítulo~\ref{ch:cap12} usará la ELBO y el truco de reparametrización como piedras angulares
del VAE, y la KL entre distribuciones en tiempos distintos como función de pérdida
de los modelos de difusión.

%───────────────────────────────────────────────────────────────────────────────
\label{sec:estrella7}
\begin{estrella}{Cuantificación de incertidumbre en el sustituto neutrónico}

\textbf{El problema.}
El sustituto MLP $f_{\mathrm{MLP}}(\bm{u}) \approx \bm{\Phi}(\bm{u})$
(Capítulos~\ref{ch:cap4}--\ref{ch:cap6}) proporciona predicciones
deterministas del flujo neutrónico. En el contexto regulatorio de un reactor
de investigación, una predicción puntual no es suficiente: se necesita una
estimación de la \textbf{incertidumbre epistémica} sobre la predicción
---cuán seguros estamos del valor predicho dado que el modelo se entrenó con
datos finitos--- y de la \textbf{incertidumbre aleatoria} ---el ruido
intrínseco del flujo simulado por Monte Carlo.

\textbf{MC Dropout como estimador de incertidumbre operacional.}
El procedimiento estándar para cuantificar incertidumbre sin reentrenar el
modelo es el MC Dropout (Definición~\ref{def:mc_dropout}): se realizan $K$
evaluaciones del sustituto con dropout activo, produciendo $K$ predicciones
$\{\bm{\Phi}^{(k)}\}_{k=1}^K$. La media $\bar{\bm{\Phi}} = \frac{1}{K}\sum_k\bm{\Phi}^{(k)}$
es la predicción puntual; la desviación estándar por componente,
$\mathrm{std}[\Phi_j] = \sqrt{\frac{1}{K}\sum_k(\Phi_j^{(k)}-\bar{\Phi}_j)^2}$,
cuantifica la incertidumbre epistémica de forma localizada sobre el núcleo.

\textbf{Propagación de incertidumbre en secciones eficaces.}
Los parámetros nucleares del código de referencia (secciones eficaces de fisión,
captura y dispersión) tienen incertidumbres propias de $\sim 1$--$5\%$ derivadas
de los experimentos de medición nuclear. El marco bayesiano permite propagar
estas incertidumbres: si las secciones eficaces tienen prior
$p(\bm{\sigma}_x) = \mathcal{N}(\bar{\bm{\sigma}}_x, \bm{\Sigma}_{\sigma})$
con covarianza conocida $\bm{\Sigma}_\sigma$ (disponible en bibliotecas nucleares
como ENDF/B), la incertidumbre en el flujo predicho es:
\[
  \mathrm{Var}[\bm{\Phi}] \approx \bm{J}_\sigma\,\bm{\Sigma}_\sigma\,\bm{J}_\sigma^\top,
  \qquad
  \bm{J}_\sigma = \frac{\partial f_{\mathrm{MLP}}}{\partial\bm{\sigma}_x},
\]
donde el jacobiano $\bm{J}_\sigma$ se obtiene por backpropagation (Capítulo~\ref{ch:cap6})
a través del sustituto. Esto reproduce, sobre el sustituto diferenciable, el
análisis de sensibilidad que el código TSUNAMI realiza sobre el transporte
de Boltzmann.

\textbf{Calibración del sustituto.}
Un sustituto bien entrenado pero mal calibrado puede predecir flujos con
confianza excesiva en regiones del espacio de operación pobremente muestreadas.
El error de calibración esperado (ECE, Definición~\ref{def:ece}) proporciona
una métrica cuantitativa: el sustituto está bien calibrado si la incertidumbre
predicha por MC Dropout coincide con el error real contra el código de referencia.
En la práctica, el temperature scaling ---dividir las salidas del sustituto por
un escalar $T$ optimizado en un conjunto de validación--- mejora la calibración
sin cambiar la exactitud media, con un costo computacional despreciable.

\textbf{Conexión con el capítulo.}
Las tres incertidumbres de la Sección~\ref{subsec:tres_incertidumbres} aparecen
simultáneamente: la incertidumbre aleatoria es el ruido Monte Carlo de las
simulaciones de referencia; la epistémica es la incertidumbre del sustituto
sobre regiones pobremente exploradas del espacio de operación; la de modelo es
el sesgo sistemático de un MLP de capacidad finita respecto al transporte exacto.
El marco bayesiano, y en particular MC Dropout, permite separar y cuantificar
las tres de forma computacionalmente eficiente sobre el sustituto.
\end{estrella}

%───────────────────────────────────────────────────────────────────────────────
\section*{Ejercicios computacionales}
\addcontentsline{toc}{section}{Ejercicios computacionales}
%───────────────────────────────────────────────────────────────────────────────

\begin{ejerciciocomp}{Ajuste probabilístico: MSE vs.\  log-verosimilitud de Poisson}
\textbf{Objetivo:} verificar que la pérdida óptima depende del modelo de ruido
subyacente, y que usar MSE para conteos de Poisson produce estimaciones sesgadas.

\medskip\noindent\textbf{Algoritmo:}
\begin{enumerate}
  \item Generar $N = 200$ muestras de un proceso de Poisson:
        $\lambda(x) = 5\exp(-x^2/2)$ con $x_i \sim \mathcal{U}[-3,3]$ y
        $y_i \sim \mathrm{Poisson}(\lambda(x_i))$.
  \item Ajustar un MLP $(1,32,1)$ con activación softplus (para garantizar
        $\hat{\lambda} > 0$) usando dos pérdidas: MSE y log-verosimilitud de
        Poisson $\ell = -\sum_i [y_i\log\hat{\lambda}(x_i) - \hat{\lambda}(x_i)]$.
  \item Para cada ajuste, graficar $\hat{\lambda}(x)$ frente a $\lambda(x)$
        real. ¿En qué región difieren más? ¿Cuál produce predicciones negativas?
  \item Calcular el estadístico $\chi^2$ reducido para cada ajuste. ¿Cuál
        está mejor calibrado?
\end{enumerate}

\medskip\noindent\textbf{Verificación:} para conteos bajos ($y_i \in \{0,1,2\}$),
el MLE de Poisson debe ser significativamente mejor que MSE; para conteos altos
($\bar{y} \gg 1$) las dos pérdidas deben converger.
\end{ejerciciocomp}

\begin{ejerciciocomp}{MC Dropout para cuantificación de incertidumbre}
\textbf{Objetivo:} implementar MC Dropout e interpretar la incertidumbre predicha
en términos de las regiones exploradas vs.\ no exploradas por los datos de
entrenamiento.

\medskip\noindent\textbf{Algoritmo:}
\begin{enumerate}
  \item Entrenar un MLP $(1,64,64,1)$ con dropout ($p_d = 0.1$) y pérdida MSE
        sobre $N=50$ puntos de entrenamiento con $x \in [-3,3]$ muestreados de
        $f(x) = \sin(x) + \varepsilon$, $\varepsilon \sim \mathcal{N}(0,0.1^2)$.
  \item Con dropout activo en test, hacer $K=100$ evaluaciones sobre el rango
        $x \in [-5,5]$. Calcular la media $\bar{y}(x)$ y el intervalo
        $[\bar{y}-2\mathrm{std}, \bar{y}+2\mathrm{std}]$.
  \item Graficar la banda de incertidumbre. ¿Crece la incertidumbre fuera del
        rango de entrenamiento? ¿Es este comportamiento consistente con la
        distinción entre incertidumbre aleatoria y epistémica?
  \item Comparar con un MLP idéntico sin dropout (predicción determinista).
        ¿Qué información adicional provee MC Dropout?
\end{enumerate}
\end{ejerciciocomp}

%───────────────────────────────────────────────────────────────────────────────
\begin{notasbib}
La interpretación probabilística del aprendizaje supervisado tiene raíces en la
estadística matemática clásica. El principio de máxima verosimilitud fue
formalizado por Fisher (1922)~\cite{fisher1922mathematical}, que demostró
su consistencia y eficiencia asintótica. La cota de Cramér-Rao fue establecida
independientemente por Cramér (1946) y Rao (1945), y formalizada en la forma
de la información de Fisher. El tratamiento moderno puede encontrarse en
Lehmann y Casella (1998)~\cite{lehmann1998theory}.

El principio de máxima entropía fue formulado por Jaynes (1957)~\cite{jaynes1957information}
como generalización del principio de equiprobabilidad de Laplace: en ausencia de
información adicional, el prior correcto es el que maximiza la entropía sujeto
a las restricciones disponibles. Jaynes argumentó que la mecánica estadística
de Gibbs-Boltzmann es un caso especial de inferencia bayesiana bajo restricciones
de energía. Esta perspectiva conecta directamente con la distribución estacionaria
del SGD del Capítulo~\ref{ch:cap5}.

La conexión entre entropía cruzada y divergencia KL, y la derivación de la
pérdida de entropía cruzada como MLE bajo modelo categórico, son resultados
estándar de la teoría de la información (Cover y Thomas, 2006)~\cite{cover2006elements}.
El cuello de botella de información como principio organizador del aprendizaje
fue propuesto por Tishby, Pereira y Bialek (2000)~\cite{tishby2000information}.

La inferencia variacional como técnica de aproximación al posterior fue
desarrollada en la comunidad de estadística bayesiana computacional. El truco
de reparametrización aplicado a redes neuronales y la derivación del estimador
SGVB son debidos a Kingma y Welling (2013)~\cite{kingma2013auto}, que lo
introdujeron en el contexto del autoencóder variacional (VAE). La
interpretación del dropout como inferencia variacional aproximada fue establecida
por Gal y Ghahramani (2016)~\cite{gal2016dropout}.

La teoría VC fue desarrollada por Vapnik y Chervonenkis (1971)~\cite{vapnik1971uniform}
como fundamento matemático del aprendizaje estadístico. El tratamiento moderno
con análisis del margen y cotas de norma espectral puede encontrarse en Bartlett
et al.\ (2017)~\cite{bartlett2017spectrally}. Las limitaciones de las cotas
de VC en el régimen sobredeterminado moderno y el fenómeno \textit{double descent}
se discuten en Belkin et al.\ (2019)~\cite{belkin2019reconciling} y Bartlett
et al.\ (2020)~\cite{bartlett2020benign}.

El gradiente natural fue introducido por Amari (1998)~\cite{amari1998natural}
en el contexto de la geometría de la información, y la conexión entre Adam y el
gradiente natural diagonal fue señalada por varios autores. El algoritmo K-FAC
(Kronecker-factored Approximate Curvature), que extiende esta idea a bloques de
matrices, fue propuesto por Martens y Grosse (2015)~\cite{martens2015optimizing}.

En el contexto de la física de reactores, la propagación de incertidumbre en las
secciones eficaces y el análisis de sensibilidad mediante el método adjunto son
implementados en el módulo TSUNAMI del paquete SCALE (Oak Ridge National
Laboratory). La biblioteca de datos nucleares ENDF/B-VIII.0 proporciona las
matrices de covarianza $\bm{\Sigma}_\sigma$ que permiten el análisis bayesiano
de incertidumbre descrito en la sección estrella.
\end{notasbib}

%───────────────────────────────────────────────────────────────────────────────
\section*{Ejercicios resueltos}
\addcontentsline{toc}{section}{Ejercicios resueltos}
%───────────────────────────────────────────────────────────────────────────────

\begin{ejercicio}
\label{ej:predictor_optimo}
\textbf{(La esperanza condicional como predictor óptimo).}

Demostrar que $f^*(\bm{x}) = \mathbb{E}[y|\bm{x}]$ minimiza el error cuadrático
medio poblacional $R(f) = \mathbb{E}_{p_{\text{data}}}[(f(\bm{x})-y)^2]$ sobre
todas las funciones medibles $f$.

\medskip\noindent\textit{Solución.}
Añadiendo y restando $f^*(\bm{x})$ dentro del cuadrado:
\begin{align*}
  R(f) &= \mathbb{E}\bigl[(f(\bm{x})-y)^2\bigr] \\
  &= \mathbb{E}\bigl[(f(\bm{x})-f^*(\bm{x}))^2\bigr]
    + 2\,\mathbb{E}\bigl[(f(\bm{x})-f^*(\bm{x}))(f^*(\bm{x})-y)\bigr]
    + \mathbb{E}\bigl[(f^*(\bm{x})-y)^2\bigr].
\end{align*}
El término cruzado es cero: condicionando sobre $\bm{x}$,
$\mathbb{E}[(f^*(\bm{x})-y)|\bm{x}] = f^*(\bm{x}) - \mathbb{E}[y|\bm{x}] = 0$.
Por tanto:
\[
  R(f) = \underbrace{\mathbb{E}[(f(\bm{x})-f^*(\bm{x}))^2]}_{\geq 0}
  + \underbrace{\mathbb{E}[(f^*(\bm{x})-y)^2]}_{=\,\mathrm{Var}[y|\bm{x}]}.
\]
El primer término es no negativo y se anula si y solo si $f = f^*$. El segundo
es la varianza del ruido irreducible, independiente de $f$. Luego $f^*$ minimiza
$R(f)$. $\square$
\end{ejercicio}

\begin{ejercicio}
\label{ej:dropout_vi}
\textbf{(Dropout como inferencia variacional: caso escalar).}

Sea una BNN con un único peso $w$ con prior $p(w) = \mathcal{N}(0,1)$ y familia
variacional $q_\phi(w) = p_d\,\delta(w) + (1-p_d)\,\mathcal{N}(w;m,\sigma^2)$
con $p_d = 0.5$ y $\sigma^2 \to 0$. Calcular la KL $D_{\text{KL}}(q_\phi\|p)$
y mostrar que la ELBO se reduce a la pérdida con dropout más regularización.

\medskip\noindent\textit{Solución.}
En el límite $\sigma^2 \to 0$, $q_\phi \approx 0.5\,\delta(w) + 0.5\,\delta(w-m)$.
La KL entre una mezcla de dos puntos y una gaussiana es:
\begin{align*}
  D_{\text{KL}}(q_\phi\|p)
  &= 0.5\,\log\frac{q_\phi(0)}{p(0)}
   + 0.5\,\log\frac{q_\phi(m)}{p(m)} \\
  &= 0.5\log\frac{0.5/\sigma}{p(0)}
   + 0.5\log\frac{0.5/\sigma}{p(m)}.
\end{align*}
El término que depende de $m$ es $-0.5\log p(m) = 0.5\log(2\pi) + \frac{m^2}{4}$.
La ELBO es:
\[
  \mathrm{ELBO} = \mathbb{E}_q[\log p(\mathcal{D}|w)] - D_{\text{KL}}
  \approx 0.5\log p(\mathcal{D}|0) + 0.5\log p(\mathcal{D}|m)
  - \frac{m^2}{4} + \mathrm{const.}
\]
El término $0.5\log p(\mathcal{D}|0) + 0.5\log p(\mathcal{D}|m)$ es la
log-verosimilitud esperada bajo dropout con $p_d = 0.5$: con probabilidad $0.5$
el peso es cero (neurona inactiva) y con probabilidad $0.5$ es $m$ (neurona
activa). El término $-m^2/4$ es regularización $\ell_2$ sobre el peso. $\square$
\end{ejercicio}

\begin{ejercicio}
\label{ej:kl_gaussianas}
\textbf{(Divergencia KL entre gaussianas).}

Calcular $D_{\text{KL}}(\mathcal{N}(\mu_1,\sigma_1^2)\|\mathcal{N}(\mu_2,\sigma_2^2))$
en forma cerrada.

\medskip\noindent\textit{Solución.}
\begin{align*}
  D_{\text{KL}}(p\|q)
  &= \int p(x)\log\frac{p(x)}{q(x)}\,dx \\
  &= \mathbb{E}_p\!\left[\log p(x) - \log q(x)\right] \\
  &= \mathbb{E}_p\!\left[
    -\frac{(x-\mu_1)^2}{2\sigma_1^2} + \frac{(x-\mu_2)^2}{2\sigma_2^2}
    + \log\frac{\sigma_2}{\sigma_1}
  \right].
\end{align*}
Usando $\mathbb{E}_p[(x-\mu_1)^2] = \sigma_1^2$ y $\mathbb{E}_p[(x-\mu_2)^2]
= \sigma_1^2 + (\mu_1-\mu_2)^2$:
\[
  D_{\text{KL}}(\mathcal{N}(\mu_1,\sigma_1^2)\|\mathcal{N}(\mu_2,\sigma_2^2))
  = \log\frac{\sigma_2}{\sigma_1}
  + \frac{\sigma_1^2 + (\mu_1-\mu_2)^2}{2\sigma_2^2} - \frac{1}{2}.
\]
Este resultado se usará en el VAE del Capítulo~\ref{ch:cap12}, donde la KL entre el posterior
variacional $q_\phi(\bm{z}|\bm{x}) = \mathcal{N}(\bm{\mu}_\phi, \mathrm{diag}
(\bm{\sigma}_\phi^2))$ y el prior $p(\bm{z}) = \mathcal{N}(\bm{0},\bm{I})$
tiene la misma forma con $\mu_2 = 0$ y $\sigma_2 = 1$. $\square$
\end{ejercicio}

%───────────────────────────────────────────────────────────────────────────────
\section*{Ejercicios propuestos}
\addcontentsline{toc}{section}{Ejercicios propuestos}
%───────────────────────────────────────────────────────────────────────────────

\begin{ejercicio}
\textbf{(MLE para distribuciones distintas).}
\begin{enumerate}[label=(\alph*)]
  \item Para $p(y|\bm{x};\bm{\theta}) = \mathrm{Laplace}(f(\bm{x};\bm{\theta}),b)$
        (distribución de Laplace con media $f$ y escala $b$), mostrar que el MLE
        es el minimizador de la pérdida MAE (\textit{mean absolute error}):
        $\mathcal{L}_{\text{MAE}} = \frac{1}{N}\sum_i|f(\bm{x}_i;\bm{\theta})-y_i|$.
  \item Para $p(y|\bm{x};\bm{\theta}) = \mathrm{Poisson}(f(\bm{x};\bm{\theta}))$
        con $f > 0$, calcular el MLE y la forma resultante de la pérdida.
  \item ¿Qué distribución generativa corresponde a la pérdida de Huber, que
        combina MSE para errores pequeños y MAE para errores grandes?
\end{enumerate}
\end{ejercicio}

\begin{ejercicio}
\textbf{(Principio de máxima entropía).}
\begin{enumerate}[label=(\alph*)]
  \item Usando multiplicadores de Lagrange, derivar que la distribución de
        máxima entropía sobre $\mathbb{R}$ sujeta a $\mathbb{E}[X] = \mu$ y
        $\mathbb{E}[X^2] = \mu^2+\sigma^2$ es $\mathcal{N}(\mu,\sigma^2)$.
  \item Derivar que la distribución de máxima entropía sobre $\mathbb{R}_+$
        sujeta a $\mathbb{E}[X] = \mu$ es la distribución exponencial
        $p(x) = \frac{1}{\mu}e^{-x/\mu}$.
  \item ¿Qué restricción corresponde al prior de Jeffreys para un modelo
        gaussiano de media desconocida y varianza fija?
\end{enumerate}
\end{ejercicio}

\begin{ejercicio}
\textbf{(La ELBO para un modelo gaussiano simple).}
Sea $p(\bm{z}) = \mathcal{N}(\bm{0},\bm{I})$, $p(\bm{x}|\bm{z}) =
\mathcal{N}(\bm{W}\bm{z},\sigma^2\bm{I})$ (un modelo lineal gaussiano) y
$q_\phi(\bm{z}|\bm{x}) = \mathcal{N}(\bm{\mu}_\phi(\bm{x}),
\bm{\Sigma}_\phi(\bm{x}))$ con $\bm{\Sigma}_\phi = \mathrm{diag}(\bm{s}_\phi^2)$.
\begin{enumerate}[label=(\alph*)]
  \item Calcular la ELBO en forma cerrada usando el resultado del Ejercicio
        Resuelto~\ref{ej:kl_gaussianas}.
  \item Implementar el truco de reparametrización para este modelo y verificar
        que el gradiente de la ELBO respecto a $\bm{\mu}_\phi$ y $\bm{s}_\phi$
        puede calcularse con backpropagation.
  \item ¿Cuál es la distribución variacional óptima $q^*(\bm{z}|\bm{x})$ para
        este modelo? ¿Coincide con el posterior exacto?
\end{enumerate}
\end{ejercicio}

\begin{ejercicio}
\textbf{(Calibración y temperatura).}
Sea un clasificador binario con logit $z(\bm{x}) = f(\bm{x};\hat{\bm{\theta}})$
y predicción $\hat{p}(\bm{x}) = \sigma(z(\bm{x}))$.
\begin{enumerate}[label=(\alph*)]
  \item En un conjunto de validación de $N=1000$ muestras, calcular el ECE
        para 10 bins de igual tamaño con las probabilidades predichas y las
        etiquetas reales (puede usarse un dataset sintético).
  \item Aplicar temperature scaling: $\hat{p}_T(\bm{x}) = \sigma(z(\bm{x})/T)$
        y optimizar $T$ para minimizar la pérdida de entropía cruzada en el
        conjunto de validación.
  \item ¿Mejora el ECE con la temperatura óptima? ¿Cambia la precisión del
        clasificador?
  \item Repetir con MC Dropout ($K=50$ evaluaciones). ¿Es el modelo con
        dropout mejor calibrado que sin él?
\end{enumerate}
\end{ejercicio}

\begin{ejercicio}
\textbf{(Información de Fisher para el MLP).}
Sea un MLP con una capa oculta y pérdida de entropía cruzada.
\begin{enumerate}[label=(\alph*)]
  \item Calcular el elemento diagonal $I_{jj}(\bm{\theta}) =
        \mathbb{E}[(\partial\log p/\partial\theta_j)^2]$ para un peso de la
        capa de salida.
  \item Comparar $I_{jj}$ con el segundo momento del gradiente $\hat{v}_j$
        de Adam después de muchas iteraciones. ¿Son proporcionales?
  \item El gradiente natural para este modelo sería
        $\bm{I}^{-1}\nabla\mathcal{L}$. ¿Cuál es el costo computacional de
        calcular $\bm{I}^{-1}$ exactamente vs.\ la aproximación diagonal de
        Adam?
  \item Para la aproximación de Kronecker-factored (K-FAC, Martens y Grosse,
        2015), $\bm{I} \approx \bm{A} \otimes \bm{B}$ donde $\bm{A}$ y $\bm{B}$
        son matrices de menor dimensión. ¿Cuánto se reduce el costo de inversión?
\end{enumerate}
\end{ejercicio}


\printbibliography[heading=subbibliography]